\chapter{Réglages déjà intégrés}
\label{chap:reglages-existants}

\section{Convention de calcul}

Le pipeline manipule un pixel $\mathbf{p}=(r,g,b)^T$ en valeurs flottantes. Les calculs photométriques doivent être réalisés sur des composantes linéaires ; le codage sRGB ou le format d'export est appliqué à la fin. La luminance relative utilisée pour les normalisations est :
\begin{equation}
  Y = 0{,}2627r + 0{,}6780g + 0{,}0593b.
\end{equation}

La formule est celle du master Rec.~2020 linéaire D65, définie dans \texttt{WorkingColorSpace.h}, et non celle du signal sRGB de sortie. Conserver une précision flottante avant l'encodage réduit les pertes de détails dans les hautes lumières et évite les saturations intermédiaires inutiles.

\section{Tonalité globale}

\subsection{Exposition, contraste et tons}

L'exposition en valeurs IL (EV) est un gain multiplicatif :
\begin{equation}
  \mathbf{p}' = 2^{E}\mathbf{p},
\end{equation}
où $E$ est le réglage d'exposition. Une variation de $+1$~EV double donc l'énergie des composantes linéaires.

Le contraste et les réglages \emph{hautes lumières}, \emph{ombres}, \emph{blancs} et \emph{noirs} sont pilotés par des fonctions de pondération de luminance. De façon générique, une correction locale de tonalité prend la forme :
\begin{equation}
  Y' = Y + a\,w(Y),
\end{equation}
où $a$ est l'intensité demandée et $w$ une fenêtre douce concentrée sur la plage ciblée. Les composantes sont ensuite remises à l'échelle avec conservation de chromaticité. Cette approche évite qu'un réglage d'ombres affecte autant les hautes lumières.

Cette équation décrit le principe, pas l'algorithme complet de \texttt{ToneRecovery}, qui exploite aussi des statistiques de l'image source. Les validations sRGB historiques ne qualifient pas automatiquement le master Rec.~2020. La matrice, le rapport évolutif Lumières et les JSON de \texttt{docs/validation} précisent les campagnes et les limites ; une validation sur corpus n'est pas une garantie pour tous les boîtiers.

\subsection{Courbes}

Les courbes disponibles sont une courbe de luminance et des courbes Rouge, Vert et Bleu. Entre les points de contrôle $(x_i,y_i)$, l'implémentation s'appuie sur une interpolation polynomiale de Lagrange :
\begin{equation}
  P(x)=\sum_{i=0}^{n} y_i\prod_{j\ne i}\frac{x-x_j}{x_i-x_j}.
\end{equation}

Les courbes de canal permettent une teinte créative, tandis que la courbe de luminance règle le contraste tonal. Les valeurs HDR sont prolongées en linéaire au-delà de l'intervalle de contrôle afin de ne pas introduire de coupure brutale.

\section{Balance des blancs et teinte}

La température est calculée à partir du lieu planckien corrélé selon l'algorithme de Krystek~\cite{krystek1985}. Les gains de chromaticité sont appliqués au moyen d'une adaptation chromatique de Bradford~\cite{lam1985}, formulée matriciellement :
\begin{equation}
 \mathbf{RGB}' = M_{XYZ\rightarrow RGB}\,B^{-1}\,D\,B\,M_{RGB\rightarrow XYZ}\,\mathbf{RGB},
\end{equation}
où $B$ est la matrice de Bradford et $D$ une matrice diagonale de rapports entre blanc de référence et blanc cible dans l'espace LMS.

La teinte compense l'axe vert--magenta. Les gains colorés sont normalisés pour conserver autant que possible la luminance :
\begin{equation}
  n=0{,}2627r_g+0{,}6780g_g+0{,}0593b_g, \qquad \mathbf{p}'=\frac{\mathbf{g}\odot\mathbf{p}}{n}.
\end{equation}
Cette écriture illustre la normalisation de gains dans le master ; le chemin RAW comprend également les matrices du profil et la référence de blanc source. Les essais du 2 octobre 2026 sont historiques : consulter le rapport Couleur et la matrice pour la qualification du chemin courant.

\section{Couleur globale, mélangeur et color grading}

\subsection{Vibrance et saturation}

La saturation ajuste uniformément l'amplitude chromatique. La vibrance protège davantage les couleurs déjà saturées. Si $c=(\max(r,g,b)-\min(r,g,b))/\max(r,g,b)$ est la chroma relative, une approximation du gain de vibrance est :
\begin{equation}
  g_v = 1+\frac{V}{100}(1-c).
\end{equation}
Les couleurs peu saturées reçoivent donc l'essentiel de l'effet. Cette description correspond à \texttt{ColorAdjustments} dans \texttt{dev} au 22 octobre 2026 ; les variantes expérimentales étudiées sur \texttt{patch\_color} ne sont pas présumées intégrées. Les vérifications historiques ne constituent pas une validation de toutes les variantes.

\subsection{Mélangeur de couleur}

Le module C++ \texttt{ColorMixer} agit dans une représentation HSV du RGB linéaire
sur huit bandes : rouge, orange, jaune, vert, cyan, bleu, violet et magenta.
Les centres, en degrés, sont :
\[0,\;30,\;60,\;120,\;180,\;240,\;270,\;300.\]
Pour une teinte d'origine $h$, l'influence d'une bande centrée en $h_i$ est triangulaire :
\begin{equation}
  w_i(h)=\max\left(0,1-\frac{|\delta_i(h)|}{a_i(h)}\right),
\end{equation}
où $\delta_i(h)$ est la distance circulaire signée dans $[-180^\circ,180^\circ]$ et $a_i(h)$ la distance du centre au centre voisin du côté considéré. Les poids sont calculés sur la teinte source et multipliés par $q=t^2(3-2t)$, avec $t=\operatorname{clamp}(S/0{,}05,0,1)$, pour protéger les achromatiques. Les commandes pondérées sont additionnées avant application, sans gains multipliés aux chevauchements.
\begin{align}
 H'&=H+18^\circ\sum_i qw_iH_i/100,\\
 S'&=\operatorname{clamp}\left(S\left[1+\sum_i qw_iS_i/100\right],0,1\right),\\
 Y'&=Y\,2^{0{,}5\sum_i qw_iL_i/100}.
\end{align}
La reconversion HSV est normalisée à la luminance source avant le gain sélectif. Les bornes $\pm100$ garantissent un gain entre $1/\sqrt{2}$ et $\sqrt{2}$. Le seuil de chroma $0{,}05$ et la borne $0{,}5$ EV sont des choix initiaux empiriques, pas une calibration Lightroom validée. Les gris sont inchangés exactement et les canaux négatifs ne sont pas écrêtés dans ce module. Le gamut mapping de sortie peut modifier $Y$. Le mélangeur et le grading restent indépendants de la saturation et de la vibrance globales ; un mélangeur à zéro ne change pas le rendu de base.

\subsection{Color grading}

Trois zones indépendantes sont disponibles : ombres, tons moyens et hautes lumières. Elles reçoivent une couleur définie par teinte et saturation. Une fonction lisse de $Y$ distribue l'influence de chaque zone ; la correction chromatique est ajoutée avec une intensité proportionnelle à la saturation, puis bornée au moment de l'export. Cette solution est volontairement légère et non destructive ; une future validation visuelle doit préciser ses limites en saturation forte.

\section{Détail local, voile et grain}

\subsection{Texture et clarté}

Texture et clarté reposent sur la composante de détail d'une décomposition par flou gaussien :
\begin{equation}
  D_\sigma = Y-G_\sigma(Y), \qquad Y'=Y+aD_\sigma w(Y),
\end{equation}
où $G_\sigma$ est un flou gaussien. Texture utilise une échelle courte ($\sigma\approx1{,}2$), clarté une échelle plus large ($\sigma\approx8$). Les poids de protection limitent l'accentuation dans les extrêmes et les zones très peu chromatiques. La primitive de flou est cohérente avec celle documentée par OpenCV~\cite{opencvGaussian}.

\subsection{Correction de voile}

Le modèle physique simplifié de Koschmieder relie l'observation $I$, la radiance de scène $J$, la lumière atmosphérique $A$ et la transmission $t$ :
\begin{equation}
  I(x)=J(x)t(x)+A(1-t(x)).
\end{equation}
L'estimation par \emph{dark channel prior} utilise localement $\min_c I_c$ puis reconstruit $J$ après estimation et bornage de $t$~\cite{he2009,he2011}. L'implémentation applique cette logique comme une correction graduelle afin d'éviter les halos.

\subsection{Réduction de bruit et netteté CPU}

\paragraph{Modèle introduit le 5 octobre 2026, porté au master Rec.~2020.}

Le débruitage précède les effets locaux et la netteté. Il sépare le RGB
linéaire en $Y=0{,}2627R+0{,}6780G+0{,}0593B$, $Co=R-B$ et
$Cg=G-(R+B)/2$. Les plans ont une estimation de bruit distincte par le MAD
du détail diagonal Haar $(p_{00}-p_{01}-p_{10}+p_{11})/2$, normalisé par
$0{,}67448975$. Les blocs exactement constants ne participent pas à cette
estimation ; celle-ci reste un proxy après dématriçage, pas un modèle capteur.

Pour $s=\operatorname{clamp}(v,0,100)/100$, les résidus de deux flous
successifs ($\sigma=0{,}8$ puis $1{,}8$ pixels) subissent un seuil
$\lambda_f=\widehat\sigma ms$, avec $m=2{,}55$ pour Y et $3{,}15$ pour
Co/Cg ; $\lambda_m=0{,}52\lambda_f$. Un guide préfiltré et un Sobel
normalisé protègent les bords, y compris les frontières chromatiques
isoluminantes. Le résultat est mélangé à l'original avec le poids $s$.
Les changements positifs et négatifs sont équilibrés pour conserver la
moyenne sans ajouter un offset aux noirs. Ces coefficients sont empiriques.

La reconstruction s'écrit $RGB=Y+d$, où la luminance de $d$ est nulle.
Si un canal devient négatif, réduire $d$ vers le gris de même Y, sans
clamp supérieur HDR. Le réglage couleur seul préserve ainsi la luminance
linéaire par pixel à la précision flottante. À bruit 0/0, le chemin est
court-circuité exactement. La conservation de moyenne globale n'interdit
pas une dérive locale ; une recette photographique reste nécessaire.

La netteté CPU est un masque flou sur Y, avec seuil relatif au bruit et
limitation locale des dépassements. Son estimateur historique et sa
calibration ne changent pas dans cette correction. Le test natif
\texttt{detailProcessingTests} couvre les régressions synthétiques, pas une
parité Lightroom ni une fidélité garantie de toutes les textures.

\subsection{Grain}

Le grain est un bruit pseudo-aléatoire déterministe, dépendant de la position du pixel et de l'intensité du réglage. Sa déterminisation assure qu'un même export produit le même résultat et évite le scintillement entre deux rendus successifs.

\section{Persistance et réglage par lot}

L'état de présentation ne fait pas partie du sidecar photographique.
Le fichier global \texttt{ipr-ui.ini} mémorise les sections ouvertes, les
sélections de canal/bande/zone, les couleurs principales, la visibilité et les
dimensions des panneaux ainsi que le format et la qualité d'export. La
restauration est protégée contre les écritures de valeurs par défaut QML au
démarrage ; elle ne modifie pas les réglages de développement.
Les clés réservées au recadrage, aux masques, à l'IA et au GPU sont inactives.
Le menu Affichage donne accès à la visibilité des panneaux et au chemin de
l'INI, qui se personnalise application fermée.

Le sidecar utilise une signature \texttt{IPR1}, un numéro de version, la taille et la date de modification du fichier source. Il stocke les réglages, les courbes et les métadonnées utilisateur via \texttt{QDataStream}~\cite{qtQDataStream}. L'écriture différée de 450~ms passe par \texttt{QSaveFile}, qui écrit temporairement puis valide le fichier de façon atomique~\cite{qtQSaveFile}. Une empreinte source non concordante invalide le sidecar plutôt que d'appliquer potentiellement des réglages à un autre fichier.

Le format courant est v6 (\texttt{Sidecar.h}), et non v4. Il conserve les huit triplets du mélangeur, les choix de développement RAW, les références et empreintes de ressources ainsi que les réglages optiques. Les lectures d'anciens schémas présentes dans le code ne garantissent pas un rendu identique après changement de modèle. Le numéro de schéma de l'INI est indépendant de celui du sidecar.

L'application par lot copie seulement les réglages de développement de l'image active vers les images cochées. Les annotations locales de destination sont préservées lorsqu'elles existent. Ce choix évite qu'une correction photographique n'écrase involontairement un classement déjà établi.
